\section{Experimental Setup}
\label{sec:experiments}

\paragraph{Models.} Qwen3-8B \citep{qwen3} and Llama-3.1-8B \citep{llama31}, both dense 8B-scale decoder-only transformers under permissive licences. Qwen3 is the April 2025 release of Alibaba's open-weight family; we use the non-thinking instruction-tuned variant for deterministic evaluation. We deliberately exclude MoE releases (Llama~4 Scout/Maverick, Qwen3-MoE) because our causal-tracing methodology targets dense attention and MLP projections; extending \fvs{} to routed experts is future work. Cross-architecture-family transferability (e.g., Mistral, Phi-3) is outside the scope of this study and the claim is correspondingly scoped (\S\ref{sec:limitations}).

\paragraph{Task sequence.} Four publicly available English tasks, chosen for diversity of output structure: (i) NER on CoNLL-2003 \citep{sang2003conll}; (ii) extractive QA on SQuAD-v1.1 \citep{rajpurkar2016squad}; (iii) abstractive summarisation on XSum \citep{narayan2018xsum}; (iv) code generation, training on CodeAlpaca-20k \citep{codealpaca} and evaluated on HumanEval \citep{chen2021humaneval}. To test order sensitivity, we additionally run a \textit{reversed} order (Code $\rightarrow$ Summarisation $\rightarrow$ QA $\rightarrow$ NER) for the Qwen3-8B configuration and compare \fvs{} rankings (\S\ref{sec:results:ablations}).

\paragraph{Metrics.} Per-task: span-F1 (NER, \texttt{seqeval}), exact-match / F1 (QA), ROUGE-L (summarisation), pass@1 (code, sandboxed execution). We report average prior-task forgetting $\bar{\mathrm{Forget}}^{(K)}$ and backward-transfer (BWT) \citep{lopez2017gem}. All tables report mean $\pm$ standard deviation over 3 seeds. Statistical significance on forgetting comparisons is assessed with paired $t$-tests across seeds; we also report Wilcoxon signed-rank $p$-values when the normality assumption is violated.

\paragraph{Baselines.}
(1) Full fine-tuning; (2) Standard LoRA ($r{=}16$) on all attention and MLP projections; (3) EWC (online Fisher) \citep{kirkpatrick2017ewc}; (4) SDFT \citep{shi2024sdft}; (5) LoRA + episodic replay; (6) DoRA \citep{liu2024dora}; (7) VeRA \citep{kopiczko2024vera}; (8) LOFIT \citep{yin2024lofit}, which localises LoRA targets by \textit{new-task adaptation} rather than prior-task resistance and is thus the closest counterpoint to \tglora{}; (9) a random-target control that applies LoRA to a uniformly-random $50\%$ subset of $\mathcal{A}_{\mathrm{std}}$ matched in parameter count. All baselines share base models, task order, learning-rate schedule, batch size, and random seeds.

\paragraph{Hyperparameters.} AdamW \citep{loshchilov2019adamw} with $\beta_1{=}0.9$, $\beta_2{=}0.999$, weight decay $0.01$. Learning rate $2\!\times\!10^{-4}$ for LoRA-style methods and $2\!\times\!10^{-5}$ for full fine-tuning, selected over $\{1, 2, 5\}\!\times\!10^{-5}$ for full-FT and $\{1, 2, 5\}\!\times\!10^{-4}$ for LoRA by validation-forgetting on a held-out slice of stage-0 data. Per-device batch size $4$, gradient accumulation $4$ (effective batch $16$), $3$ epochs per stage, warmup ratio $0.03$ ($\approx100$ steps at our scale), max sequence length $1024$ (raised to $2048$ for summarisation). bf16 mixed precision with gradient checkpointing. LoRA $r{=}16$, $\alpha{=}32$, dropout $0.05$.

\paragraph{Compute.} All experiments run on NVIDIA A40 48GB GPUs. Seeds: $\{42, 1337, 2024\}$. Main matrix: $2\text{ models} \times 9\text{ methods} \times 3\text{ seeds} = 54$ sequential-FT runs, each covering 4 task stages, for $\approx \pend{total_a40_hrs}$ A40-hours.
